\documentclass[11pt]{article}
\usepackage[a4paper,margin=1.85cm]{geometry}
\usepackage[T1]{fontenc}
\usepackage{lmodern}
\usepackage{amsmath,amssymb,amsthm,mathtools}
\usepackage{graphicx}
\usepackage{xcolor}
\usepackage{microtype}
\usepackage{enumitem}
\usepackage[hidelinks]{hyperref}

\definecolor{ink}{HTML}{172033}
\definecolor{accent}{HTML}{385B7A}
\definecolor{soft}{HTML}{EEF5FA}
\color{ink}
\newtheorem{theorem}{Theorem}
\newtheorem{lemma}{Lemma}
\newtheorem{remark}{Remark}

\newcommand{\D}{\mathcal D}
\newcommand{\F}{\mathsf F}
\newcommand{\Hil}{\mathcal H}

\title{The commutator subgroup of Thompson's group $F$ is framenable}
\author{Solution packet for Problem 5 in arXiv:2508.20061}
\date{11 August 2026}

\begin{document}
\maketitle

\begin{center}
\fcolorbox{accent}{soft}{\parbox{0.91\textwidth}{
\textbf{Full positive answer.}
Let $F'=[F,F]$ and let $\D=\mathbb Z[1/2]\cap(0,1)$.  The permutation representation of $F'$ on $\ell^2(\D)$ weakly contains the trivial representation, because every finite collection of elements fixes all dyadic points sufficiently close to $0$.  Meanwhile, choosing one group element for each image of a fixed point turns a sparse orbit into the standard orthonormal basis.  Thus the same representation has almost invariant vectors and a weak frame indexed by a countably infinite subset, exactly as required by framenability.}}
\end{center}

\section*{Source question}
Dutkay, Georgescu, and Picioroaga introduce framenability in Definition 3.1 of \cite{Source}.  A countably infinite discrete group $G$ is framenable when there are a unitary representation $\pi:G\to\mathcal U(\Hil)$, a countably infinite subset $S\subset G$, and $v\in\Hil$ such that
\[
 1_G\prec\pi
 \quad\text{and}\quad
 \{\pi_s v:s\in S\}\text{ is a weak frame in }\Hil.
\]
Here a weak frame is a total Bessel family.  In Remark 4.11(v), on source PDF page 21, they write:
\begin{quote}
``Our tools do not capture the commutator $[F,F]$ as framenable. We think this should be the case.''
\end{quote}

\begin{center}
\includegraphics[width=0.92\textwidth]{figures/open_problem_crop.png}
\end{center}

\section*{Proof intuition}
The definition permits the frame to be indexed by a subset of the group.  This is decisive for a permutation representation.  If $G$ acts transitively on an infinite set $X$, fix $x_0\in X$ and retain exactly one group element sending $x_0$ to each point of $X$.  The corresponding orbit vectors of $\delta_{x_0}$ are then the orthonormal basis of $\ell^2(X)$; stabilizer repetitions have disappeared.

For $F'$, the action on the interior dyadic points is transitive.  It also has ``fixed points escaping to the endpoint'': every $g\in F'$ is the identity near $0$, so $g(2^{-n})=2^{-n}$ for all large $n$.  The delta masses at these points are almost invariant unit vectors.  These two features occur in the same permutation representation and give framenability immediately.

\section*{A transitive-action criterion}
\begin{lemma}[escaping fixed points]\label{lem:criterion}
Let $G$ be a countably infinite group acting transitively on a countably infinite set $X$.  Suppose there is a sequence $(x_n)$ in $X$ such that
\begin{equation}\label{eq:eventual}
 \text{for every }g\in G,\qquad gx_n=x_n\quad\text{for all sufficiently large }n.
\end{equation}
Then $G$ is framenable.  In fact, its witnessing weak frame can be chosen to be an orthonormal basis.
\end{lemma}

\begin{proof}
Let $\pi$ be the permutation representation on $\Hil=\ell^2(X)$,
\[
 \pi_g\delta_x=\delta_{gx}\qquad(g\in G,\ x\in X).
\]
Fix $x_0\in X$.  Transitivity makes the orbit map $q:G\to X$, $q(g)=gx_0$, surjective.  Choose a section of $q$, equivalently a subset $S\subset G$ on which $q$ is bijective.  Since $X$ is countably infinite, so is $S$, and
\[
 \{\pi_s\delta_{x_0}:s\in S\}=\{\delta_x:x\in X\}
\]
is the standard orthonormal basis of $\ell^2(X)$.  In particular, it is a Parseval frame and hence a weak frame.

Put $\xi_n=\delta_{x_n}$.  These are unit vectors, and \eqref{eq:eventual} gives
\[
 \|\pi_g\xi_n-\xi_n\|=0
 \qquad\text{for every fixed }g\in G\text{ and all sufficiently large }n.
\]
Thus $(\xi_n)$ is a sequence of almost invariant unit vectors, so $1_G\prec\pi$.  The representation $\pi$, vector $\delta_{x_0}$, and subset $S$ satisfy the definition of framenability.
\end{proof}

\section*{Application to Thompson's commutator subgroup}
Recall that Thompson's group $F$ consists of the orientation-preserving piecewise-linear homeomorphisms of $[0,1]$ whose breakpoints are dyadic and whose slopes are integral powers of $2$.  Write
\[
 F'=[F,F],\qquad \D=\mathbb Z[1/2]\cap(0,1).
\]
We use two standard facts \cite[Theorem 4.1 and Lemma 4.2]{CFP}:
\begin{enumerate}[leftmargin=2em]
\item $F'$ consists exactly of the elements of $F$ that are the identity on neighborhoods of both $0$ and $1$;
\item $F'$ acts transitively on $\D$.
\end{enumerate}

For completeness, the second fact can also be seen from the usual transitivity of $F$ on $\D$.  Given $x,y\in\D$, start with $f\in F$ satisfying $f(x)=y$.  The endpoint-slope homomorphism
\[
 \rho:F\longrightarrow\mathbb Z^2,
 \qquad
 \rho(f)=\bigl(\log_2 f'(0+),\log_2 f'(1-)\bigr)
\]
has kernel $F'$.  Copies of $F$ supported on $[0,y]$ and $[y,1]$ fix $y$ and allow the two endpoint slopes to be corrected independently.  Hence there is $u\in\operatorname{Stab}_F(y)$ with $\rho(u)=-\rho(f)$; then $u\circ f\in F'$ and $(u\circ f)(x)=y$.

\begin{theorem}\label{thm:main}
The commutator subgroup $[F,F]$ of Thompson's group $F$ is framenable.
\end{theorem}

\begin{proof}
Let $F'$ act on $\D$ in the natural way.  It is a countable group acting transitively on the countably infinite set $\D$.  Set $x_n=2^{-n}$ for $n\geq1$.  For any fixed $g\in F'$, the endpoint-triviality theorem gives $\varepsilon_g>0$ such that $g(t)=t$ on $[0,\varepsilon_g]$.  Therefore
\[
 g(x_n)=x_n\qquad\text{whenever }2^{-n}<\varepsilon_g.
\]
Thus the escaping-fixed-point condition \eqref{eq:eventual} holds.  Lemma~\ref{lem:criterion} proves that $F'$ is framenable.
\end{proof}

Concretely, the witness in Theorem~\ref{thm:main} is the representation
\[
 \pi:F'\to\mathcal U(\ell^2(\D)),\qquad
 \pi_g\delta_x=\delta_{g(x)}.
\]
For any $x_0\in\D$, choose a transversal $S$ for the stabilizer of $x_0$.  Then $\{\pi_s\delta_{x_0}:s\in S\}$ is an orthonormal basis, while $(\delta_{2^{-n}})$ is an almost invariant sequence.

\section*{Why the source's orbit obstruction does not apply}
Remark 4.10 of \cite{Source} observes that the same type of permutation representation does not have a weak frame vector when the orbit is indexed by the \emph{whole} group: infinitely many elements in a point stabilizer repeat the same nonzero coefficient, violating the Bessel inequality.  There is no contradiction.  Framenability was deliberately defined using a countably infinite subset $S\subset G$.  The transversal above contains one representative of each orbit point, so there are no repetitions and the selected orbit is an orthonormal basis.

The result proves only framenability.  It does not produce a weak frame indexed by all of $F'$ and does not settle the amenability of $F$ or $F'$.

\section*{Verification and bounded novelty check}
The proof was audited directly against Definition 3.1: $S$ is countably infinite, the selected orbit is Parseval, and the almost invariant vectors belong to the same representation.  The standard facts about $F'$ were cross-checked against \cite{CFP} and later sources that quote the same theorem and lemma.

On 11 August 2026, the run indexes and bounded web/arXiv searches were checked for the source identifier, ``framenable'', ``commutator subgroup'', ``Thompson'', and direct variants of ``$[F,F]$ is framenable''.  The latest located source version still states the question, and no later paper giving this answer was found.  The 2026 journal landing page reported no citing papers.  This is evidence for novelty, not a definitive guarantee.

\section*{Human-review recommendation}
Review as a likely valid full positive solution to Problem 5.  The only imported group-theoretic input is classical and explicit; once transitivity and endpoint triviality of $F'$ are confirmed, the representation-theoretic argument is immediate.

\begin{thebibliography}{9}
\bibitem{Source}
D. E. Dutkay, C. Georgescu, and G. Picioroaga,
\emph{Frame Vector Group Representations and Amenability Properties},
arXiv:2508.20061v3; J. Math. Anal. Appl. (2026), article 130837.

\bibitem{CFP}
J. W. Cannon, W. J. Floyd, and W. R. Parry,
\emph{Introductory notes on Richard Thompson's groups},
Enseign. Math. (2) \textbf{42} (1996), 215--256.
\end{thebibliography}

\end{document}
